\section{一维环因子扫描与有理逼近层级}

\subsection{环旋转动力学}

由 $(\alpha,\mathcal{P})$ 定义环相位 $\theta_t\in S^1$ 的推进
\[
\theta_{t+1}=\theta_t+\alpha\pmod 1.
\]

\subsection{有理/无理与周期/稠密}

线性流理论表明：在环面上线性流轨道在有理相关时周期、在有理无关时稠密。限制到 $S^1$ 因子即得：$\alpha\in\QQ$ 时轨道周期闭合，$\alpha\notin\QQ$ 时轨道稠密（参见 \cite{WikipediaSturmianWord} 的相关讨论与引用路径）。
有限时间与有限精度下对无理 $\alpha$ 的可达实现可由有理逼近序列 $p_k/q_k$ 表达。

\subsection{与六维 strip 近似体的一致性}

六维超立方模型的结论给出几何层的对应：有理 $(1/1)$ strip 的投影产生周期结构，无理 $(\tau/1)$ strip 的投影产生准晶结构 \cite{PRB46_6091}。
本文将其作为“有理逼近—无理极限”在几何层与符号层的对照。

\subsection{有理逼近—周期轨道—有限样本统计（桥接引理）}

\begin{lemma}[有理参数的周期性]\label{lem:rational-periodic}
若 $\alpha=p/q\in\QQ$（既约），则环旋转轨道 $\theta_{t+1}=\theta_t+\alpha\ (\mathrm{mod}\ 1)$ 的符号序列 $w_t=\mathbf{1}_{I_1}(\theta_t)$ 为周期序列，其周期整除 $q$。
\end{lemma}

\begin{proposition}[有限样本的近似层级（可检验命题）]\label{prop:finite-sample-approx-hierarchy}
设 $\alpha\notin\QQ$，取其连分数收敛子序列 $\alpha_k=p_k/q_k\to\alpha$。对任意固定检验口径（例如长度 $n$ 的因子频率、复杂度估计或平衡性偏差），用 $\alpha_k$ 生成的周期序列在足够大的但有限观测窗内可作为对无理极限统计的近似层级。该命题在本文仅作为实验可检验的桥接口径使用，不在一般性上断言统一收敛速率。
\end{proposition}

